% To produce Postscript and PDF:
%    latex template; latex template; 
%    dvips -o template.ps template; ps2pdf template.ps
\documentclass[a4paper,11pt]{article}
\usepackage{mathptmx}
%\usepackage{newtxmath}   not available ? 
%\usepackage{newtxtext} []  not available ?
\usepackage[margin=40mm]{geometry}% <-------- CHANGE HERE for the global margins
\usepackage[T1]{fontenc}    % Check that ÖÄÅöäå come out ok!
\usepackage[utf8]{inputenc}
\usepackage{graphicx}
\usepackage{url} 
\usepackage{parskip}
\usepackage{listings} % for code snippets
% should not USE  ? 
%\setlength{\parindent}{0mm} % Do not indent the 1st line of a paragraph.
%\setlength{\parskip}{3mm}   % Add space between paragraphs.

% defines my code style
\usepackage{color}
\usepackage{courier}

\definecolor{dkgreen}{rgb}{0,0.6,0}
\definecolor{gray}{rgb}{0.5,0.5,0.5}
\definecolor{mauve}{rgb}{0.58,0,0.82}

\lstset{frame=tb,
  language=Java,
  aboveskip=3mm,
  belowskip=3mm,
  showstringspaces=false,
  columns=flexible,
  basicstyle={\scriptsize\ttfamily},
  numbers=left,
  numberstyle=\tiny\color{gray},
  keywordstyle=\color{blue},
  commentstyle=\color{dkgreen},
  stringstyle=\color{mauve},
  breaklines=true,
  breakatwhitespace=true,
  tabsize=2
}
 

% Save some paper by stuffing more text on each page:
% A4: 210mm x 297mm, approximately 35 mm margins on every side.
\addtolength{\topmargin}{-14mm}    
\addtolength{\textheight}{24mm}    
\addtolength{\oddsidemargin}{-4mm} 
\addtolength{\textwidth}{14mm}     

\renewenvironment{abstract}
{\small
  \begin{center}
  \bfseries \abstractname\vspace{-.5em}\vspace{0pt}
  \end{center}
  \list{}{
    \setlength{\leftmargin}{1cm}%
    \setlength{\rightmargin}{\leftmargin}%
  }%
  \item\relax}
{\endlist}

% This creates the header.
\makeatletter
\renewcommand{\@oddhead}
{\fontsize{10}{12}\selectfont \hfill OpenGL ES 3.0 and WebGL 1.0 on Android platform \hfill }
\makeatother

\begin{document}

%============================================================
% title
\label{Title} 
\title{OpenGL ES 3.0 and WebGL 1.0 on Android platform}
\author{Patryk Małek\\
        University of Novi Sad\\
        Faculty of Sciences\\
        \texttt{malekpatryk@gmail.com}
        }
\date{}         % Do not print the date on the final paper!
\maketitle
% prevents page numbering on this page
\thispagestyle{empty}

%============================================================

\vspace{7pc}

\centerline{
\includegraphics[width=0.3\textwidth,height=0.3\textheight,keepaspectratio]{NoviSadLogoGray.jpg}
}

\vspace{9pc}

\begin{abstract}
\label{Abstract}
Android OS has gained much popularity in the past few years, spreading from mobile devices to tv sets, automotive industry and multimedia tablets. 
To meet users' growing expectations of multimedia content delivery methods and their demand for new astonishing features new open standards, like HTML5 and WebGL, were introduced. 
The past few years have been very fruitful for developers working 2D and 3D visualisation techiniques for embedded systems and many new framworks and APIs have been developed.
\newline The goal of this paper is to present basic concepts, history of OpenGL ES 3.0 and WebGL 1.0 along with exemplar code snippets for those frameworks/APIs.
Presented code will be targeted for Android platform.
\end{abstract}
\pagebreak

%============================================================

%\tableofcontents

%============================================================

\section{Introduction} 
\subsection{Background}
Nowadays, mobile devices such as cell phones or PDAs have become the most ubiquitous ones of all, especially Android platform which is rapidly gaining popularity.
In the last months of 2013 Android reached 81\% market share in smartphone market.
With increasing demand for new features, user experiences and ongoing evolution of chips,  developers and designers are working hard to keep up with users' desires using state of the art hardware and sofware.
Nowadays, the high performance graphic processors of spartphones and tablets easily handle 3Dvisualization, and with no difficulty regarding shortage of memory.

In the past few years there has been a lot of development in the field of 3D visualisation in embedded systems e.g.\ Android platform which has given life to many frameworks or APIs like OpenGL ES, WebGL, three.js and many others.

\subsection{OpenGL ES}
\emph{OpenGL for Embedded Systems} \cite{opengles_kronos} is a subset of \emph{OpenGL} \cite{opengl_kronos} API for rendering 2D and 3D graphics created by \emph{Kronos Group} \cite{kronos_group}.
It has been designed to be used with embedded systems like smartphones, tablets or video consoles.
OpenGL ES shows a good example of re-constructing a general purpose desktop 3D graphics library to a small, low-level rendering library for embedded systems.
OpenGL ES aims to provide an extremely compact API without sacrificing features.
Its primary objective (for the release of version 1.0) was to be fully implementable in software in under 50kB of code while being well-suited for hardware acceleration \cite{mobile_3d_graphics_with_OGLES_M3G}.
Nowadays mobile devices are able to support graphics effects similar to the desktop versions, thanks to the the library that is compressed in size, and still packed with features.
\newline OpenGL ES always provides compatibility with OpenGL's version that it has been based on.
This way it is allowing developers to immediately port their mobile version of the application to desktops and then add desktop specific features or effects to their product.
This flexibility of OpenGL ES allows developers for faster and easier porting of applications between mobile and desktop devices.

First version of OpenGL ES that has been released was marked as version 1.0 and it was based on API from OpenGL 1.3.
It kept most of the functionality from OpenGL and added some as well (for instance, removed need for wraping OpenGL calls with \emph{glBegin} and \emph{glEnd}).
\newline In 2004 Kronos Group OpenGL 1.1 which improved image quality and optimizations to increase performance while reducing memory bandwidth usage to save power.
\newline In March 2007 OpenGL ES 2.0 has been released and enabled fully programmable 3D graphics (thanks to programmable pipeline).
It was based on OpenGL 2.0 and this compatibility lasted until the release of OpenGL 4.1.
\newline The most recent version of OpenGL ES was released in August 2013 and it is marked as 3.0 \cite{opengl_es3_spec}.
It has been based on OpenGL 4.3 (but one might say that respectively it can be placed somewhere between OpenGL 3.1 and 4.3).
Version 3.0 is also backwards compatible with OpenGL ES 2.0 enabling developers to add new video processing features later on in the development process when they decide to support newer version of the API.

\subsection{WebGL}
\emph{WebGL} \cite{webgl_kronos} is a cross platform, web standard JavaScript API for rendering low-level 3D content in any compatible web browser without use of any external plugins like Adobe Flash.
It can be also reffered to as Implementation of OpenGL ES for the web.
WebGL functionality is exposed through HTML5's Canvas element as DOM (Document Object Model) interface.
Nowadays most of the current versions of web browsers (including the mobile versions) support WebGL.
WebGL grew out of the Canvas 3D experiments started by Vladimir Vukićević at Mozilla and has been presented in 2006 as a Canvas 3D prototype.
By the end of 2007, both Mozilla and Opera had made their own separate implementations.
In early 2009, the non-profit technology consortium Khronos Group started the WebGL Working Group, with initial participation from Apple, Google, Mozilla, Opera, and others.
Version 1.0 of the WebGL specification was released in March 2011.
Early applications of WebGL include Google Maps and Zygote Body \cite{zygote_body}.
More recently, Autodesk has ported most of their applications to the cloud running on local WebGL clients.
The initial version of WebGL 2.0 specification has been released in September 2013 and its specification is based on OpenGL ES 3.0.

\subsection{Android platform}
\emph{Android} \cite{androidcom} is an operating system (based on Linux operating system \cite{gnulinux}), primarily designed for smartphones but after time it has also been developed for tablets, tvs etc.
The very first unofficial versions of Android came out in late 2008 but the first official, to be supported by Google came out on 30th April 2009 marked as version 1.5 - \emph{Cupcake}.
\newline Currently it is the most popular operating system for mobile devices with around 80\% of market share.
OpenGL ES and Android already have quite a history together.
First version of OpenGL ES API marked as 1.0 was implemented in Android platform in version 1.0 marked as - \emph{Apple Pie}.
After that there has been more improvements and features implemented with each version.
Since version 4.3 released in 2013, Android will support OpenGL ES in version 3.0.

\subsection{Deployment platform - Google Nexus 4}
For deployment purposes \emph{Google Nexus 4} has been chosen.
Google branded LG smarthone, released on 13th November 2013.
\newline Having mid-high end specs: processor Qualcomm Snapdragon™ S4 Pro with 4 cores each 1.5 Ghz and 2GB of RAM it will definately fit purposes as a deployment platform.
\newline With its recent update from Google to Android 4.3 it has received support for OpenGL ES 3.0 together with many software improvements like: enhancements for rendering pipeline, a new version of GLSL\footnote{OpenGL Shading Language is a high-level shading language based on syntax of the C programming language, created by OpenGL ARB (OpenGL Architecture Review Board) to give developers more direct control over graphics pipeline without having to use ARB assembly language or hardware-specific languages.} shading language or greatly enhanced texturing functionalities.

%============================================================

\pagebreak[3]
\section{OpenGL ES}

Sec. \ref{sec:math} gives examples of math formulas.

% reference to section examplae
% (Sec.~\ref{sec:emphasis})
%
%------------------------------------------------------------

\pagebreak[3] 
\subsection{Mathematics}\label{sec:math}

Longer formulas may be put on a separate
line:
\[ \emptyset \in \emptyset \; \Rightarrow \; E \neq mc^2. \]

You may also want to number the formulas like Eqn.~\ref{eqn:myequation1}
below.
\begin{equation}\label{eqn:myequation1}
C = E_{K_{public}}(P) = P^e. \hspace{10mm}   P = D_{K_{private}}(C) = C^d.
\end{equation}

%------------------------------------------------------------

\subsection{Make a list}\label{sec:list}

Lists can have either bullets or numbers on them. 

\begin{itemize}
\item one item
\item another item, which is an exceptionally long one for an item
  and consequently continues on the next line.
\end{itemize}

Lists can have several levels. Item~\ref{kukkuu} below contains
another list.
\begin{enumerate}
\item the fist item \label{kukkuu}
\item the second item
\end{enumerate}

%============================================================

\clearpage 
\section{Programming Examples} 

\lstinputlisting[caption={Exemplar piece of code in Java using OpenGL ES API}]{example1.java}

%============================================================

\section{Summary} 

%============================================================

\clearpage
\label{Bibliography} 
%\pagebreak[4]
%
\bibliographystyle{plain}
\bibliography{bibliography} 
%
% In that case, remember to run bibtex:
% latex template; bibtex template; latex template; latex template; 

\end{document}